% Placeholders. Red text marks anything that is not yet a result or a decision.
% grep for \placeholder before any submission; check-arxiv-llm-compliance flags leftovers.
\newcommand{\placeholder}[1]{\textcolor{red}{[#1]}}

% \placeholderfig{width}{height}{text}: a framed box that stands in for a figure.
\newcommand{\placeholderfig}[3]{%
  \fbox{\parbox[c][#2][c]{\dimexpr#1-2\fboxsep-2\fboxrule\relax}{\centering\textcolor{red}{\small #3}}}}

% Names used often.
\newcommand{\IC}{\textsc{InferCell}}
\newcommand{\syn}{JCVI-syn3A}
\newcommand{\Mzero}{M0}
\newcommand{\Rhat}{\widehat{R}}
\newcommand{\sigb}{\sigma_{\mathrm{b}}}
\newcommand{\code}[1]{\texttt{#1}}

% /write-paper markers. \todo{} is red text for a result or conclusion that does not exist yet
% (M10, KIM-116, KIM-90). \stub{} is a drafting note, \DRAFT{} a number not yet sourced.
% check_paper.py --final catches \stub and \DRAFT but NOT \todo: before submission, grep for \todo{
% (check-arxiv-llm-compliance flags the red text). check_prose.py counts \todo text as prose.
\newcommand{\todo}[1]{\textcolor{red}{\textbf{TODO:} #1}}
\newcommand{\stub}[1]{\textcolor{red}{[stub: #1]}}
\newcommand{\DRAFT}[1]{\textcolor{red}{#1}}
